\section{Robust responses and certified dense recovery}\label{sec:robust-screening}

The response description admits two different extensions. First, the
leader can protect itself against followers whose cost is close to the
true optimum, without requiring those followers to be stationary.
Second, a structured surrogate can identify enough coordinate statuses
to recover the exact response of a dense convex quadratic follower.
The first extension changes the set of responses considered by the leader;
the second preserves the original follower and certifies its solution.

\subsection{Robustness through criterion-specific measurements}\label{subsec:robust-model}

Use the moving-normal quadratic-block model of
\cref{subsec:moving-normals}, including its compact domain $C$,
polynomial coordinate bounds, and positive-definite local quadratic
matrices. The fixed-normal model is included. The follower value $v(x)$
and cost-near-optimal set $S_\rho(x)$ are those of
\eqref{eq:follower-value} and \eqref{eq:nearoptimal-set}, with a supplied polynomial
budget $\rho(x)\ge0$ on $C$. Let each upper criterion, including the
objective indexed by zero, have the form
\begin{equation}\label{eq:measured-criterion}
 G_j(x,z)=p_j(x,T_j(x)z),\qquad j=0,\ldots,m,
\end{equation}
where $T_j(x)$ has $h_j\le h$ rows and $p_j,T_j$ are explicitly encoded
rational polynomials. Only $h$ is fixed; the number of criteria and the
rank of their combined measurement matrices may grow. An affine upper
row needs only one measurement, regardless of its follower support.
Leader-only terms belong directly to $p_j$. Affine measurements with an
additional offset can be written in this form by absorbing the offset
into $p_j$.

The robust feasible leaders and objective are
\begin{align}
 D_\rho&=\{x\in C:P(x)\ne\varnothing,\
      G_j(x,z)\le0\text{ for every }z\in S_\rho(x),\ j=1,\ldots,m\},
 \label{eq:robust-domain}\\
 W_\rho(x)&=\max_{z\in S_\rho(x)}G_0(x,z),\qquad
                    \operatorname{val}_\rho=\inf_{x\in D_\rho}W_\rho(x).
 \label{eq:robust-objective}
\end{align}
The budget is measured from the global follower minimum. Upper rows
never remove responses from $S_\rho(x)$ before an adversarial choice.
The near-optimality model and separate upper-row adversaries are established
in \cite{BesanconAnjosBrotcorne2024}; related complexity-class membership
results and an objective-robust variant appear in
\cite{BesanconAnjosBrotcorne2021}. The following theorem
concerns exact global optimization for the stated structural dimensions.

\begin{samepage}
\begin{theorem}[Exact near-optimal robustness]\label{thm:nearoptimal-robust}
Fix $(r,d,s,k_=,k_{\le},h)$. For
\eqref{eq:robust-domain}--\eqref{eq:robust-objective}, feasibility, the
finite infimum when feasible, and attainment are computable in polynomial
bit time in $(L,\delta)$. If attained, an optimal leader and a worst
near-optimal response can be returned jointly in one real-algebraic field
with polynomial total encoding length. At any supplied rational leader
with $P(x)\ne\varnothing$, the exact worst value and a witnessing
near-optimal response for any specified criterion can also be computed
in polynomial bit time. For many separately requested criteria, the
polynomial encoding guarantee applies separately to each witness;
a common field containing all unrelated witnesses is not promised.
\end{theorem}
\end{samepage}

\begin{lemma}[Measurement-fiber threshold]\label{lem:measurement-threshold}
Fix one measurement matrix $T(x)$ with at most $h$ rows. Let
\[
 P_T(x,\tau)=\{z\in P(x):T(x)z=\tau\},\qquad
 m_T(x,\tau)=\min_{z\in P_T(x,\tau)} f(x,z)
\]
when the fiber is nonempty. At a follower-feasible leader,
\begin{equation}\label{eq:measurement-threshold}
 \tau\in T(x)S_\rho(x)\quad\Longleftrightarrow\quad
 P_T(x,\tau)\ne\varnothing\ \text{ and }\
                       m_T(x,\tau)\le v(x)+\rho(x).
\end{equation}
Moreover the right side is equivalent to existence of a feasible KKT
candidate of the measurement-fiber problem whose cost is at most
$v(x)+\rho(x)$. That candidate need not be a stationary point of the
unrestricted follower problem.
\end{lemma}
\begin{proof}
If a near-optimal response has measurement $\tau$, minimizing on its
compact fiber cannot increase cost. Conversely, a fiber minimizer whose
cost meets the budget is itself a near-optimal response with that
measurement. A global fiber minimizer satisfies polyhedral KKT conditions,
so it supplies the asserted candidate whenever the right side holds.
Conversely every feasible fiber candidate meeting the budget is already
a near-optimal point. Its global optimality on the fiber is unnecessary
for this converse. Global comparison remains essential in computing the
nominal value $v(x)$.
\end{proof}

\begin{proof}[Proof of \cref{thm:nearoptimal-robust}]
If $C$ is empty, report infeasibility. Otherwise first obtain the nominal value graph $\mathcal V(x,v)$ by eliminating
the response variables from the global predicate
\eqref{eq:global-response-formula}. It is quantifier-free of polynomial
size, and it has exactly one value $v=v(x)$ at a feasible leader and
none otherwise.

For one criterion $j$, add $T_j(x)z=\tau$ as $h_j$ shared equalities.
The parameters are now $(x,\tau)$, of fixed dimension $r+h_j$.
All possible measurements lie in a rational box with polynomial encoding
length: maximize the absolute values of the polynomial entries of
$T_j$ and the coordinate bounds on compact $C$, using fixed-dimensional
real algebra, and round the resulting bounds outwards to integers.
Summing the products of these bounds gives a bound for each measurement.
Univariate root bounds applied to the polynomial-size algebraic outputs
ensure the integer bounds have polynomial bit length. Thus the enlarged
parameter domain can be taken compact. The positive-definite local
quadratics are unchanged, and the added normals are polynomial, so
\cref{lem:moving-compression} applies.

Let $\mathcal E_j^{\rm fib}(x,\tau,\xi_j)$ be the quantifier-free
candidate predicate constructed as in \eqref{eq:candidate-formula},
including the new equality multipliers. Its coordinate $t_j$ is the
candidate's actual follower cost. Define
\begin{equation}\label{eq:criterion-attainable-formula}
 \mathcal A_j(x,v,q)\ \Longleftrightarrow\
 \exists\tau,\xi_j\,
 [\mathcal E_j^{\rm fib}(x,\tau,\xi_j)\wedge
       t_j\le v+\rho(x)\wedge q=p_j(x,\tau)].
\end{equation}
By \cref{lem:measurement-threshold}, at $v=v(x)$ this describes
exactly all attainable criterion values on $S_\rho(x)$. A representative
may be a minimizer on its measurement fiber even when the response it
replaces is nonstationary for the unrestricted follower.

Eliminate the variables in \eqref{eq:criterion-attainable-formula}
\emph{separately for every criterion}, producing quantifier-free
$\mathcal A_j$. Then eliminate $q$ in
\[
             \mathcal B_j(x,v)\ \Longleftrightarrow\
                 \exists q\,[\mathcal A_j(x,v,q)\wedge q>0].
\]
Each elimination has a fixed total number of variables. There are
polynomially many such computations, whose output formulas have
polynomial total length and degree. After this elimination, the formula
\begin{equation}\label{eq:robust-domain-formula}
 \mathcal D(x)\ \Longleftrightarrow\
 x\in C\ \wedge\ \exists v\,
       [\mathcal V(x,v)\wedge\bigwedge_{j=1}^m\neg\mathcal B_j(x,v)]
\end{equation}
uses only one shared nominal-value variable. The empty-follower case
is excluded by $\mathcal V$. Introducing simultaneous quantified
adversaries for all rows would not give the claimed bound.

Since $S_\rho(x)$ is nonempty and compact at every feasible leader,
its continuous criterion attains a maximum. The worst-objective graph is
\begin{equation}\label{eq:robust-worst-formula}
 \begin{split}
 \mathcal W(x,q)\ \Longleftrightarrow\quad&\mathcal D(x)\wedge
 \exists v\,[\mathcal V(x,v)\wedge\mathcal A_0(x,v,q)\ \wedge\\
 &\hspace{42mm}\neg\exists q'\,
              (\mathcal A_0(x,v,q')\wedge q'>q)].
 \end{split}
\end{equation}
Eliminate $\mathcal D$ first if desired; every formula still has a
fixed number of variables. The attained worst-value set is bounded by
uniform follower bounds and compact $C$. Its emptiness, exact infimum,
and endpoint membership follow by the construction
\eqref{eq:infimum-predicate}. Endpoint membership decides attainment.

When attained, sample the leader, nominal value and worst value together
with \emph{one} measurement and a thresholded fiber candidate realizing
that worst value, using the uneliminated relation
\eqref{eq:criterion-attainable-formula} for $j=0$. This adds only a
fixed number of variables. Simultaneous sampling returns a common field
of polynomial degree and encoding length. The rational block formulas
of the fiber candidate recover all follower coordinates in that field.
Their cost meets the true budget and their criterion equals the exact
worst value. The same procedure at a supplied rational follower-feasible
leader computes one criterion maximum and witness, including a violating
witness whenever that maximum is positive. Repeating it for different
criteria gives separate encodings and does not require a joint field
for all those witnesses.
\end{proof}

A bounded budget can be an additional leader coordinate. For example,
choosing $G_0(x,z)=-\rho$ maximizes tolerated follower cost loss under
all robust upper rows. This adds one fixed leader dimension. Exact
rational adversarial output is generally impossible: minimizing $z^2/2$
on $[0,1]$ with budget $1/4$ gives near-optimal interval
$[0,1/\sqrt2]$, whose worst response for criterion $z$ is irrational.

\begin{corollary}[Constant local matrices and algebraic degree]
\label{cor:robust-degree}
If the local $Q_b,E_b,G_b$ are independent of the leader, the common-field
outputs of \cref{thm:nearoptimal-robust} have degree bounded solely by
input degree and the fixed structural dimensions. This includes
polynomially moving shared normals and measurement matrices $T_j(x)$.
For fixed input degree, the bound is independent of the follower and
criterion counts and coefficient bit lengths. It applies to one jointly
recovered leader and adverse witness, or to each separately requested
criterion witness. It does not assert attainment.
\end{corollary}
\begin{proof}
Shared normals and added measurement rows enter the effective linear
cost and shared consistency equations; they do not enter the local KKT
matrix \eqref{eq:block-kkt-system}. That matrix is constant under the
stated assumptions, so each local response is polynomial in the
compressed variables of degree $O(\delta)$ with rational coefficients.
All candidate, threshold and criterion formulas have degree bounded
in terms of $\delta$, independent of the numbers of rows and blocks.
The degree-independent-of-polynomial-count elimination and sampling
bounds used in \cref{cor:constant-hessian-degree} apply to each separate
criterion computation and then to \eqref{eq:robust-worst-formula} and
one witness. The number of successive elimination and sampling steps along each
dependency chain is fixed. Rational
decoding preserves the resulting field.
\end{proof}

\subsection{When the robust optimum is attained}\label{subsec:robust-attainment}

\begin{proposition}[Convex fixed-normal attainment]\label{prop:robust-attainment}
Under the hypotheses of \cref{thm:nearoptimal-robust}, suppose all
follower constraint normals are fixed and $\phi(x,\cdot)$ is convex
for each $x\in C$. Then every feasible robust problem attains its
minimum, including when the nonnegative budget vanishes at some leaders.
\end{proposition}
\begin{proof}
Let $C_F=\{x\in C:P(x)\ne\varnothing\}$. It is compact: any sequence
of feasible leaders has, along a convergent subsequence, bounded feasible
followers whose limit remains feasible. Fixed normals and
\eqref{eq:hoffman} let every $z\in P(x)$ be approximated by feasible
$z_n\in P(x_n)$ whenever $x_n\to x$ in $C_F$. Closed constraints
and uniform bounds give the converse limit property.

The full follower objective is strictly convex in $z$, so it has one
nominal minimizer $z^*(x)$. The comparison-and-repair argument in
\cref{lem:closed-response}, together with uniqueness, implies continuity
of $z^*$ on $C_F$, and hence of $v(x)=f(x,z^*(x))$.
The graph of $S_\rho$ is closed by continuity of $f,v,\rho$ and the
constraints, and its values are uniformly bounded.

To prove the needed approximation property for $S_\rho$, fix
$x_n\to x$ in $C_F$ and $z\in S_\rho(x)$. If $\rho(x)=0$, then
$z=z^*(x)$ and the nearby nominal minimizers supply near-optimal
approximations. If $\rho(x)>0$, for any $\theta\in(0,1)$ the point
$z_\theta=(1-\theta)z+\theta z^*(x)$ is feasible and convexity gives
\[
 f(x,z_\theta)\le(1-\theta)f(x,z)+\theta v(x)
             \le v(x)+(1-\theta)\rho(x)<v(x)+\rho(x).
\]
Repair $z_\theta$ into $P(x_n)$ by Hoffman. Continuity of the cost,
nominal value and budget makes the repaired point near-optimal for all
sufficiently large $n$. Taking $\theta$ successively to zero and using
a diagonal sequence gives $z_n\in S_\rho(x_n)$ with $z_n\to z$.

For any continuous criterion, the upper limit of its near-optimal
maximum is bounded by its maximum at $x$: choose maximizing responses
and use a convergent subsequence and the closed graph. The reverse
inequality follows by approximating a maximizing response at $x$ using
the preceding construction. Thus each criterion maximum is continuous
on $C_F$. The robust upper rows define a closed subset of compact
$C_F$, and the continuous worst objective attains its minimum there.
\end{proof}

\begin{example}[A positive budget does not ensure attainment]
\label{ex:positive-budget-nonattainment}
Set $\rho=1/16$, $h(z)=3z^2-2z^3$, and, for $c\in[\rho,1/8]$, let
\[
                    f_c(z)=z^2(1-z)^2+ch(z),\qquad z\in[0,1].
\]
Both summands are nonnegative and $h(z)>0$ for $z>0$, so the unique
nominal optimum is zero with value zero. Separating off $z^2/2$ puts
this model in the positive-local-quadratic class with a scalar polynomial
aggregate and fixed box normals. Direct factorization gives
\[
 f_\rho(z)-\rho=(z-1)^2(z^2-z/8-1/16).
\]
Consequently $S_\rho(\rho)=[0,a]\cup\{1\}$, where
$a=(1+\sqrt{17})/16<1/3$. For $c>\rho$, the near-optimal set is
$[0,a_c]$, with $0<a_c<a$: the derivative
$f_c'(z)=2z(1-z)(1-2z+3c)$ changes sign once in $(0,1)$, and
$f_c(1)=c>\rho$. Its first crossing of the budget is therefore its
only sublevel endpoint, and $a_c\to a$ as $c\downarrow\rho$.

If the leader minimizes $c$ subject to $z\le1/2$ for every near-optimal
response, its feasible set is $(\rho,1/8]$ and its infimum $\rho$
is unattained. Even without upper rows, put $c=\rho+x$ for
$x\in[0,\rho]$ and minimize the worst value of $4x+z$. Its value
is one at $x=0$ and $4x+a_{\rho+x}$ for $x>0$. On $0<z\le a<1/3$,
\[
 f_\rho'(z)=z(4z^2-51z/8+19/8)>z/4,\qquad h(z)<z.
\]
Writing $c=\rho+x$, the identity
$x h(a_c)=f_\rho(a)-f_\rho(a_c)$ implies
$x a_c>(a-a_c)a_c/4$, and thus $4x+a_c>a$.
The objective tends to $a$ as $x\downarrow0$, so this unconstrained
robust infimum is also unattained. Unique nominal optima and a strictly
positive budget therefore do not replace the convexity assumption in
\cref{prop:robust-attainment}.
\end{example}

\begin{proposition}[An unrestricted quadratic criterion is hard]
\label{prop:robust-measurement-hardness}
Without a bound on the measurement dimension, exact worst-criterion
computation is NP-hard even for a leader-free diagonal quadratic
follower with unit-box constraints and no aggregate or resource
coupling. Universal upper-threshold feasibility is coNP-complete on
the following restricted family.
\end{proposition}
\begin{proof}
For a graph on $N$ vertices, let $f(z)=\sum_i z_i^2/2$ on $[0,1]^N$
and take budget $\rho=N/2$. The nominal value is zero, so every point
of the cube is near-optimal. Set
\[
                         G(z)=\sum_{ij\in E}(z_i-z_j)^2.
\]
This is convex. Every cube point is a convex combination of binary
vertices, so its criterion is no greater than the maximum criterion of
a binary vertex. At such a vertex, $G$ is the cut size. Unweighted
Max-Cut threshold decision is NP-complete~\cite{GareyJohnsonStockmeyer1976}.
Thus computing the worst value is NP-hard, and the robust row
$G(z)\le K-1$ is satisfied for all responses exactly when there is no
cut of size at least $K$. A violating binary vertex is a polynomial
certificate for the complement on this family, proving coNP-completeness.
This is an elementary transfer of the classical cut problem. A growing
collection of separately checked affine rows is not an arbitrary
quadratic criterion in a growing number of joint measurements.
\end{proof}

\subsection{A structured surrogate for a dense quadratic follower}
\label{subsec:screening-model}

We now consider a different model. Let $X\subset\R^r$ be a nonempty
compact rational polytope. The true follower response is
\begin{equation}\label{eq:dense-screening-follower}
 z(x)=\argmin_{0\le z\le1}
          \left\{\tfrac12 z\trans Qz+(c+Cx)\trans z\right\},
                      \qquad Q=Q\trans\succ0.
\end{equation}
All data are rational and the upper objective and rows are affine in
$(x,z(x))$. Equalities can be represented by two rows. Supply a surrogate
\[
 \widehat Q=\diag(d)+VMV\trans\succ0,\qquad
 d_i>0,\quad V\in\Q^{N\times s_0},\quad M=M\trans\in\Q^{s_0\times s_0},
\]
with the same box and linear cost. Write $y(x)$ for its exact response
and $E=Q-\widehat Q$. No norm or rank bound on $E$ is needed for
correctness. The fixed-rank construction of \cref{cor:low-rank-lp}
provides a finite closed rational polyhedral cover of the surrogate
response graph in $v=(x,w)$, $w=V\trans y(x)$. Each nonempty cover
polytope $P$ is compact, projects into $X$, and has an affine response
formula $y_P(v)$. All of its points are consistent surrogate responses,
and every leader has a graph point in the cover. The method accepts any
supplied cover with these properties, including subdivisions. Its size
and description length are part of the complexity bound.

Surrogate-based safe elimination has a direct precedent in the structured
dictionary screening of Dantas and Gribonval~\cite{DantasGribonval2019}.
Variational-inequality enclosures precede it in
\cite{LiuZhaoWangYe2014}, and complete parameter-region analysis of
QP algorithms appears in \cite{ArnstromAxehill2022}. Here the target
is an exact global upper optimum for the true dense follower, with the
exponential enumeration confined to statuses not certified on a whole
surrogate cell.

\begin{lemma}[Directional response enclosure]\label{lem:screening-enclosure}
At a fixed leader put $a=c+Cx$, $e=z-y$, $p=Ey$, and
$\widehat g=\widehat Qy+a$. Then
\begin{equation}\label{eq:screening-vi}
                  e\trans Qe+p\trans e\le0,
\end{equation}
and for any signed vector $b\in\R^N$,
\begin{equation}\label{eq:directional-enclosure}
 \left|b\trans e+\tfrac12b\trans Q^{-1}p\right|
       \le\tfrac12\sqrt{(b\trans Q^{-1}b)(p\trans Q^{-1}p)}.
\end{equation}
\end{lemma}
\begin{proof}
The two box variational inequalities are
$(Qz+a)\trans(y-z)\ge0$ and
$(\widehat Qy+a)\trans(z-y)\ge0$.
Combining the two inequalities gives \eqref{eq:screening-vi}. Completing the square yields
\[
 (e+\tfrac12Q^{-1}p)\trans Q(e+\tfrac12Q^{-1}p)
                         \le\tfrac14p\trans Q^{-1}p.
\]
Cauchy--Schwarz in the $Q$ norm gives
\eqref{eq:directional-enclosure}. This is an enclosure: points inside
the ellipsoid need not be realizable follower responses. It uses the
exact surrogate response, not an unchecked numerical approximation.
\end{proof}

\subsection{Certificates on a whole response cell}\label{subsec:whole-cell-screening}

On a nonempty cover polytope $P$, the functions $y(v)$ and $p(v)=Ey(v)$
are affine. Compute the rational number
\begin{equation}\label{eq:screening-eta}
       \eta_P=\max_{v\in\operatorname{vert}(P)}p(v)\trans Q^{-1}p(v).
\end{equation}
It is also the maximum over $P$, since the quadratic is convex and
$P$ is the convex hull of its vertices. In fixed ambient dimension
$r+s_0$, all vertices can be enumerated by rational row-basis tests,
including lower-dimensional cells and singletons, as in
\cref{app:local-support}. Define
\begin{align}
 s_i(v)&=y_i(v)-\tfrac12(Q^{-1}p(v))_i,&
 R_i&=\tfrac12\sqrt{(Q^{-1})_{ii}\eta_P},
 \label{eq:screening-centers}\\
 t_i(v)&=\widehat g_i(v)+\tfrac12p_i(v),&
 S_i&=\tfrac12\sqrt{Q_{ii}\eta_P}.
 \nonumber
\end{align}
Applying \eqref{eq:directional-enclosure} to a coordinate vector
and to its image under $Q$ gives, throughout $P$,
\begin{equation}\label{eq:screening-coordinate-gradient}
 s_i-R_i\le z_i(x)\le s_i+R_i,\qquad
 t_i-S_i\le (Qz(x)+a(x))_i\le t_i+S_i.
\end{equation}

We use labels $L,U,F$ to certify, respectively, $z_i=0$, $z_i=1$,
and zero true gradient in coordinate $i$ throughout the cell. The label
$F$ is a stationarity equation; it allows a zero-gradient bound point.
The following tests are sufficient:
\begin{center}
\begin{tabular}{cl}
\toprule
Label & Whole-cell test\\
\midrule
$L$ & $\min_P t_i>S_i$ or $\max_P s_i+R_i\le0$\\
$U$ & $\max_P t_i<-S_i$ or $\min_P s_i-R_i\ge1$\\
$F$ & $\min_P s_i>R_i$ and $\max_P s_i<1-R_i$\\
\bottomrule
\end{tabular}
\end{center}
A positive gradient in a box optimum forces the lower bound, a negative
gradient forces the upper bound, and strict interiority forces zero
gradient. Zero gradient alone does not force either bound; hence the
strict signs in the first tests for $L$ and $U$. All affine extrema
are attained at the already enumerated vertices. Square-root comparisons
are exact rational tests: for rational $A$ and $B\ge0$, for instance,
$A>\sqrt B/2$ means $A>0$ and $4A^2>B$.

The following boundary refinement avoids unnecessary ambiguity on closed
cells. The weaker conditions
\begin{equation}\label{eq:screening-closure-free}
 \min_Ps_i\ge R_i,\quad \max_Ps_i\le1-R_i,\quad
 \max_Ps_i>R_i,\quad \min_Ps_i<1-R_i
\end{equation}
also certify $F$. Indeed the affine functions $s_i-R_i$ and
$1-R_i-s_i$ are nonnegative on $P$ and neither is identically zero.
Each is positive throughout its relative interior: otherwise a small
relative-neighborhood segment through a zero interior point would force
an affine nonnegative function to vanish on the affine hull. Thus the
response is strictly interior there, giving zero gradient. The true
box-QP response is continuous in $x$. More explicitly, adding its two
variational inequalities at $x,x'$ gives
\[
 \lambda_{\min}(Q)\norm{z(x)-z(x')}_2
                               \le\norm{C(x-x')}_2
\]
when the responses differ. Zero gradient therefore extends to all of
$P$ by continuity and density of its relative interior.
Similarly, $\min_Pt_i\ge S_i$ and $\max_Pt_i>S_i$ certify $L$;
$\max_Pt_i\le-S_i$ and $\min_Pt_i<-S_i$ certify $U$.
These arguments apply in the affine hull of a lower-dimensional cell.
The strict-somewhere conditions cannot be omitted.

Choose any sound available label for a certified coordinate. Call the
others ambiguous, let $t_P$ be their number on $P$, and put
$t=\max_Pt_P$. Stronger sound certificates may reduce this number.
In particular, it measures the chosen certificates rather than an
intrinsic complexity of the response.

\begin{theorem}[Exact dense recovery after screening]\label{thm:screening-recovery}
Suppose a valid surrogate cover has $M$ nonempty rational polytopes in
fixed dimension $r+s_0$, and whole-cell sound certificates leave at most
$t$ ambiguous coordinates on each. Exact bilevel feasibility and a
rational global optimum, when feasible, can be obtained with at most
$M3^t$ rational \emph{recovery} LPs. Each LP and its preparation have
polynomial bit complexity in the data and cover descriptions. This count
excludes construction and screening of the cover. The displayed tests
provide their certificates by rational matrix computations and vertex
enumeration in fixed dimension, without additional LP calls for affine
extrema. Starting with the supplied fixed-rank surrogate, the total
bound is $3^tL^{O(r+s_0+1)}$, with $L$ including all rational input and
upper rows. No bound on the number of certified $F$ coordinates is
required.
\end{theorem}
\begin{proof}
On a cover polytope $P$, keep all certified labels and enumerate the
three labels for its ambiguous coordinates. A completed assignment
partitions the coordinates into $L,F,U$. Set $z_L=0$, $z_U=\mathbf1$,
and solve the free equations with the true Hessian:
\begin{equation}\label{eq:screening-exact-affine}
 z_F(x)=-Q_{FF}^{-1}(c_F+C_Fx+Q_{FU}\mathbf1).
\end{equation}
Every principal submatrix of an SPD matrix is SPD, so it is invertible;
the empty free set requires no inversion. Rational elimination has
polynomial bit cost even when $F$ is large and densely coupled. Impose
\begin{equation}\label{eq:screening-recovery-lp}
 \begin{gathered}
 v\in P,\qquad 0\le z_F(x)\le1,\\
 (Qz(x)+c+Cx)_L\ge0,\qquad (Qz(x)+c+Cx)_U\le0,
 \end{gathered}
\end{equation}
together with all original affine upper rows, and minimize the original
upper objective after substitution. This is a rational LP in at most
$r+s_0$ variables, bounded by compact $P$. Its feasible points satisfy
all true box KKT conditions, so they give the unique response to the
original dense follower.

Conversely, any true feasible leader--response pair has a surrogate
graph point in at least one cell. Every certified label holds there.
Assign an ambiguous coordinate label $F$ when its gradient is zero,
$L$ when its gradient is positive, and $U$ when its gradient is negative.
Box optimality makes this a compatible assignment. Thus the pair is
included in one recovery LP. Zero-gradient bound coordinates may appear
under more than one assignment, which loses no solutions. The best
nonempty LP yields an exact rational optimum; if all are empty, the
original problem is infeasible. Summing $3^{t_P}$ over the cells gives
the LP count. Fixed-dimensional cover construction and screening and
polynomial rational inversion/substitution/LP give the stated bit bound.
\end{proof}

An inexpensive additional sufficient test is often useful. Guess a full
status assignment on $P$, form \eqref{eq:screening-exact-affine}, and
check all true KKT signs and bounds at every vertex of $P$. All tested
functions are affine, so successful checks certify the response on the
entire cell without status enumeration. Failure merely means that the
assignment is not certified on the \emph{whole} cell. It must remain
available in the recovery enumeration on a smaller feasible part; a
failed guess is not an infeasibility certificate.

\subsection{A computable perturbation neighborhood}\label{subsec:screening-neighborhood}

The ambiguity bound can be certified from the geometry of the supplied
surrogate. Use its original closed clipping cells, retaining their
nominal labels $L,F,U$. At a vertex $v$, call coordinate $i$ a transition
coordinate when
\[
                    y_i(v)\in\{0,1\},\qquad\widehat g_i(v)=0.
\]
Let $q$ be the largest number of such coordinates at a vertex of any
nonempty cell. Vertices created by the leader polytope and consistency
equalities are included; no general-position assumption is made.

Every such cell has affine dimension at most $r$. Its projection to $x$
is injective, since a leader has a unique surrogate response and aggregate.
An injective linear projection on a polytope is injective on its affine
hull: a nonzero kernel direction in that affine hull would produce two
nearby points with the same projection in a relative-interior neighborhood.
Hence its dimension is at most the leader dimension.

Cover each cell by the convex hulls of all its affinely independent vertex
subsets of size at most $r+1$. These are rational simplices and they cover
the cell by the finite-dimensional convex-combination reduction proved
in \cref{lem:core-recovery}. A triangulation is unnecessary: overlapping
simplices are allowed. There are at most
$\sum_{j=1}^{r+1}\binom{V_P}{j}$ such simplices for a cell with $V_P$
vertices, polynomial in fixed dimension.

For a simplex $T$, let $J_T$ be the union of its vertices' transition
coordinates. Then $|J_T|\le(r+1)q$. If $i\notin J_T$, the containing
cell's nominal label gives strictly positive vertex margins
\[
 L:\ \widehat g_i(v),\qquad U:\ -\widehat g_i(v),\qquad
 F:\ y_i(v)\text{ and }1-y_i(v).
\]
For a nominal $F$ coordinate its surrogate gradient vanishes throughout
the closed cell, so any zero bound margin would indeed be a transition.
Let $\sigma>0$ be the minimum of all these margins over all simplices
and coordinates outside their $J_T$, setting $\sigma=1$ if the list
is empty. It is rational, and every listed margin is at least $\sigma$
throughout the corresponding simplex by affine interpolation.

For $N\ge1$, compute
\begin{equation}\label{eq:screening-radius}
 \begin{gathered}
 m_0=\frac1{\norm{\widehat Q^{-1}}_\infty},\qquad
 L_0=\norm{\widehat Q}_\infty,\qquad R=\lceil\sqrt N\rceil,\\
 \eps_0=\min\left\{\frac{m_0}{2},\
       \frac{\sigma}{4R\max\{1/m_0,\,1+L_0/m_0\}}\right\}.
 \end{gathered}
\end{equation}
Matrix infinity norms are induced maximum absolute row sums. Symmetry
implies $\norm{A}_2\le\norm{A}_\infty$, so
$\lambda_{\min}(\widehat Q)\ge m_0$ and
$\norm{\widehat Q}_2\le L_0$. The case $N=0$ is an ordinary upper LP.

\begin{corollary}[Transition-controlled dense neighborhood]
\label{cor:screening-neighborhood}
Every rational symmetric $Q$ satisfying
$\norm{Q-\widehat Q}_\infty\le\eps_0$ is positive definite. On each
of the constructed simplices, at most $(r+1)q$ coordinates require
status enumeration in \cref{thm:screening-recovery}. All radius and
cover data can be computed with polynomial bit complexity for fixed
$r,s_0$, so exact optimization is polynomial for fixed $r,s_0,q$.
The dense perturbation may have arbitrary rank and signs.
\end{corollary}
\begin{proof}
Write $\eps=\norm{E}_\infty\le m_0/2$. Then
$\lambda_{\min}(Q)\ge m_0/2$. From \eqref{eq:screening-vi},
Cauchy--Schwarz and the unit-box bound $\norm{y}_2\le R$ give
\begin{equation}\label{eq:screening-norm-response}
                          \norm{z-y}_2\le 2\eps R/m_0.
\end{equation}
The true gradient minus the surrogate gradient equals
$\widehat Q(z-y)+Ez$. Since $\norm{z}_2\le R$,
\begin{equation}\label{eq:screening-norm-gradient}
 \norm{Qz+a-\widehat g}_2
       \le 2\eps RL_0/m_0+\eps R
       \le 2\eps R(1+L_0/m_0).
\end{equation}
The choice of $\eps_0$ makes both bounds at most $\sigma/2$.
Every nontransition nominal $L$ coordinate has strictly positive true
gradient and hence is zero; a nominal $U$ coordinate has strictly
negative gradient and hence is one. A nontransition nominal $F$
coordinate lies in $[\sigma/2,1-\sigma/2]$, so its true gradient
is zero. These are sound whole-simplex certificates, independent of
whether the ellipsoid tests happened to identify them. Only $J_T$
needs enumeration.

Rational vertex enumeration, fixed-size simplex descriptions, matrix
inversion, margin evaluation and the displayed norm calculations have
polynomial bit cost. The simplex count above is polynomial for fixed
$r,s_0$; it may increase the exponent of the original arrangement
bound. Apply \cref{thm:screening-recovery} to this cover.
\end{proof}

The radius is instance-dependent and can be small. Neither low transition
multiplicity nor a numerically large margin is automatic. This is a
verifiable neighborhood of a supplied surrogate, not a general method
for finding a useful surrogate or a uniform neighborhood of every
well-conditioned Hessian.

\subsection{Value and feasibility bounds without full recovery}
\label{subsec:screening-bounds}

Even when many statuses remain ambiguous, the directional enclosure
provides computable global bounds. For a signed output vector $b$, put
\[
 q_b(v)=b\trans y(v)-\tfrac12b\trans Q^{-1}p(v),\qquad
 r_b=\tfrac12\sqrt{(b\trans Q^{-1}b)\eta_P}.
\]
Choose a rational outward upper bound $\bar r_b\ge r_b$ by rational
square-root enclosure. Its accuracy is an explicit computational choice;
any such upper bound is safe. Then
$q_b-\bar r_b\le b\trans z\le q_b+\bar r_b$ on $P$.
For a row $a_j\trans x+b_j\trans z\le h_j$, an inner sufficient
condition and an outer necessary condition are, respectively,
\begin{equation}\label{eq:screening-inner-outer}
 a_j\trans x+q_{b_j}(v)+\bar r_{b_j}\le h_j,\qquad
 a_j\trans x+q_{b_j}(v)-\bar r_{b_j}\le h_j.
\end{equation}
Use these conditions with $v\in P$ in rational inner and outer LPs,
representing equalities by both signs. If the objective is
$a_0\trans x+b_0\trans z+c_0$, minimize
$a_0\trans x+q_{b_0}-\bar r_{b_0}+c_0$ over the outer LPs.
Their best value is a certified lower bound on the true optimum:
every true feasible leader has a point in an outer LP and its objective
is no smaller than that lower enclosure. Minimize the corresponding
upper enclosure over the inner LPs for a certified upper bound and a
leader feasible for the true response. Evaluating the true follower
with an exact KKT certificate at that leader can improve this upper bound.

Empty outer LPs for every cell certify original infeasibility. Empty
inner LPs alone say nothing about original feasibility, and a minimizer
of an outer LP need not be truly feasible. The inner construction can
miss exactly feasible upper equalities. Signed objective coefficients
are handled by the directional formula, without a monotonicity
assumption. Only the actual computed finite upper-minus-lower gap gives
an objective guarantee; a small perturbation by itself supplies no
small global loss under restrictive upper rows.

\begin{example}[Small residual but complete screening ambiguity]
\label{ex:screening-conservatism}
Let $N\ge2$, $\widehat Q=I$, $Q=I+(\eps/N)\mathbf1\mathbf1\trans$,
$c=0$, $C=-\mathbf1$, and $X=[0,1]$, where $\eps>0$ is rational
and arbitrarily small. The surrogate response is $y=x\mathbf1$ and
the true response is $z=x\mathbf1/(1+\eps)$, whose gradient is zero
throughout the interval. On the whole cell $[0,1]$, however, the
coordinate enclosure radius is positive and $s_i(0)=0$, so neither
strict nor refined $F$ screening passes. The lower-gradient center
has minimum zero and positive half-width, the upper-gradient test
also fails, and the coordinate tests force neither bound. Thus the
displayed ellipsoid tests leave all $N$ coordinates ambiguous.
Here $\norm{E}_\infty=\eps$ can tend to zero. The all-$F$ guessed
assignment in \eqref{eq:screening-exact-affine} immediately certifies
the response throughout the cell and needs no enumeration. This
separates conservative screening from intrinsic difficulty.
\end{example}

Small perturbations can also move a switching point. With one follower,
$\widehat Q=1$, $Q=1+\eps$, linear cost $-2xz$, and $x\in[0,1]$,
the upper clipping point moves from $1/2$ to $(1+\eps)/2$ for
$0<\eps<1$. Reusing the surrogate's upper-bound label on its entire
cell would therefore be incorrect. Recovery must use the true KKT
conditions. Dense rational elimination and construction of many cover
cells can dominate cost even when screening is sound. The results do
not establish practical superiority, a uniformly effective cell
refinement, or automatic discovery of an accurate low-rank surrogate.
